\documentclass[11pt]{article}
\usepackage[margin=2.2cm]{geometry}
\usepackage[T1]{fontenc}
\usepackage{lmodern}
\usepackage{amsmath,amssymb}
\usepackage{booktabs}
\usepackage{parskip}
\usepackage{fancyvrb}
\usepackage{planaralgebra}
\usepackage[hidelinks]{hyperref}

\DefineVerbatimEnvironment{code}{Verbatim}
  {frame=leftline, framerule=1pt, rulecolor=\color{black!25},
   xleftmargin=1em, fontsize=\small, formatcom=\color{black!80}}
\newenvironment{result}{\par\smallskip\noindent\hspace*{1em}}{\par\medskip}
\newcommand\opt[1]{\texttt{#1}}
\newcommand\cs[1]{\texttt{\textbackslash#1}}
\newcommand\meta[1]{\ensuremath{\langle}\textit{#1}\ensuremath{\rangle}}
\newcommand\marg[1]{\texttt{\{}\meta{#1}\texttt{\}}}
\newcommand\oarg[1]{\texttt{[}\meta{#1}\texttt{]}}

\title{The \textsf{planaralgebra} package\\[2pt]
  \large Planar tangles and planar algebra diagrams in \LaTeX{} with Ti\emph{k}Z}
\date{Version 1.0, 2026/10/30}
\author{}

\PAdefine{T}{\PA{2,4}{coev, id, id}\PA{4,2}{id, ev, S(1,1)}}

\begin{document}
\maketitle

\noindent
A planar tangle is a disc with some input discs removed, together with
disjoint strings that connect the marked points on the boundaries; the
regions between the strings are shaded in a checkerboard fashion, and every
disc has a star $\star$ that marks its first region. Planar tangles act on
planar algebras by composition: $S \circ_i T$ glues the tangle $T$ into the
$i$-th input disc of $S$. This package draws such tangles from a
description layer by layer, in the style of the \textsf{plumber} package for
2d cobordisms:
\[
  \PA[numbers]{4,4}{D(2,2), A(1,1), id}\PA{4,4}{id, id, B(2,2)}\PA{4,2}{id, id, ev}
  \;\circ_3\;
  \PA{2,4}{coev, id, id}\PA{4,2}{id, ev, S(1,1)}
  \;=\;
  \PA[insert={3=T}]{4,4}{D(2,2), A(1,1), id}\PA{4,4}{id, id, B(2,2)}\PA{4,2}{id, id, ev}
\]
Each \cs{PA}\marg{m,n}\marg{items} is one layer, the side-by-side
juxtaposition of the listed strings and input discs, going from $m$ to $n$
boundary points; consecutive \cs{PA} commands are stacked into one tangle.
The package places the strings, shades the regions, draws the stars and the
outer disc, and with \opt{reduce} it computes a normal form of the tangle:
closed loops are removed and counted, and the strings are pulled tight.

\section{Quick start}

\begin{code}
\usepackage{planaralgebra}   % loads TikZ
...
\PA{2,4}{id, id, coev}\PA{4,4}{id, T(2,2), id}\PA{4,2}{id, id, ev}
\end{code}
\begin{result}
  \PA{2,4}{id, id, coev}\PA{4,4}{id, T(2,2), id}\PA{4,2}{id, id, ev}
\end{result}

The first argument is the number of boundary points above and below the
layer (\verb|2,4| means two strands in, four out); it is checked against the
items and may also be left empty. The second argument is a comma-separated
list of items, read from left to right (from top to bottom in horizontal
diagrams). \verb|id^3| is short for \verb|id, id, id|, and the same works
for any item.

\section{Generators and input discs}

\begin{center}
\begin{tabular}{@{}lllc@{}}
  \toprule
  item & aliases & strands & picture \\
  \midrule
  \opt{id}   & \opt{line}, \opt{strand} & $1\to1$ & \PA{1,1}{id} \\
  \opt{ev}   & \opt{cup}, \opt{pair}    & $2\to0$ & \PA{2,0}{ev} \\
  \opt{coev} & \opt{cap}, \opt{copair}  & $0\to2$ & \PA{0,2}{coev} \\
  \meta{label}\texttt{(}$a$\texttt{,}$b$\texttt{)} & & $a\to b$ & \PA{1,3}{x(1,3)} \\
  \bottomrule
\end{tabular}
\end{center}

Every item that is not a generator is an input disc. Its label is typeset in
math mode, so \verb|\alpha(1,1)| gives a disc labelled $\alpha$. The counts
\verb|(a,b)| say how many strands leave the disc on the input side (above
it) and on the output side (below it); \verb|(a;b)| and \verb|(a/b)| are
accepted as well. When a layer has a single disc and its counts are
declared, the counts of the disc may be omitted: \verb|\PA{2,2}{x}| is
the same as \verb|\PA{2,2}{x(2,2)}|. Discs can be repeated like
generators: \verb|f(1,1)^2|.

\cs{PAgenerators}\oarg{options} draws the items of the key \opt{list} (by
default \opt{id}, \opt{ev}, \opt{coev}, \verb|x(1,1)| and \verb|x(2,2)|) next
to each other; the starred
form leaves out \opt{id}. With \opt{names} each one gets its name.

\begin{code}
\PAgenerators[names]
\PAgenerators*[names, list={cup, cap, R(2,2), Q(1,3)}]
\end{code}
\begin{result}
  \PAgenerators[names]
\end{result}
\begin{result}
  \PAgenerators*[names, list={cup, cap, R(2,2), Q(1,3)}]
\end{result}

\section{The features}

\subsection{Composition}

Consecutive \cs{PA} commands form one tangle (spaces and single line breaks
between them are allowed; an empty line or any other token ends the
tangle). The output points of one layer are glued to the input points of the
next one, so the counts must agree. Options can be given to any \cs{PA} of
the chain and apply to the whole tangle.

\begin{code}
\[ \PA{1,3}{id, coev}\PA{3,1}{ev, id} = \PA{1,1}{id}
   \qquad
   \PA{0,2}{coev}\PA{2,2}{x}\PA{2,0}{ev} \]
\end{code}
\begin{result}
  \( \PA{1,3}{id, coev}\PA{3,1}{ev, id} = \PA{1,1}{id}
     \qquad
     \PA{0,2}{coev}\PA{2,2}{x}\PA{2,0}{ev} \)
\end{result}

If the counts do not match, or a disc has no counts that could be inferred,
the package reports an error that names the offending layer and draws a red
{\color{red}\bfseries ??} instead of the tangle. Should \cs{PA} already be
defined by another package, \cs{PAlayer} is a synonym with the same
arguments (the two can be mixed in a chain).

\subsection{Vertical or horizontal}

Tangles are \opt{vertical} by default: composition runs from top to bottom
and the first strand is on the left. With \opt{horizontal}, composition runs
from left to right. \opt{reverse} flips the direction of composition (bottom
to top, or right to left). \opt{orientation=vertical} and
\opt{orientation=horizontal} are long forms, and \cs{paset}\marg{options}
changes the defaults for the rest of the group.

\begin{code}
\PA{3,3}{id, x(2,2)}\PA{3,1}{ev, id}
\PA[horizontal]{3,3}{id, x(2,2)}\PA{3,1}{ev, id}
\PA[reverse]{3,3}{id, x(2,2)}\PA{3,1}{ev, id}
\end{code}
\begin{result}
  \PA{3,3}{id, x(2,2)}\PA{3,1}{ev, id}\qquad
  \PA[horizontal]{3,3}{id, x(2,2)}\PA{3,1}{ev, id}\qquad
  \PA[reverse]{3,3}{id, x(2,2)}\PA{3,1}{ev, id}
\end{result}

\subsection{The outer disc}

The key \opt{outer} chooses the shape of the boundary: \opt{circle} (the
default), \opt{ellipse}, \opt{rectangle} or \opt{none}. \opt{box shape}
(\opt{circle} or \opt{rectangle}) does the same for the input discs; wide
labels make rectangular discs longer.

\begin{code}
\PA[outer=ellipse]{3,3}{id, x(2,2)}\PA{3,1}{ev, id}
\PA[outer=rectangle]{3,3}{id, x(2,2)}\PA{3,1}{ev, id}
\PA[outer=none]{3,3}{id, x(2,2)}\PA{3,1}{ev, id}
\PA[box shape=rectangle]{3,1}{y(3,1)}
\end{code}
\begin{result}
  \PA[outer=ellipse]{3,3}{id, x(2,2)}\PA{3,1}{ev, id}\qquad
  \PA[outer=rectangle]{3,3}{id, x(2,2)}\PA{3,1}{ev, id}\qquad
  \PA[outer=none]{3,3}{id, x(2,2)}\PA{3,1}{ev, id}\qquad
  \PA[box shape=rectangle]{3,1}{y(3,1)}
\end{result}

\subsection{Shading and stars}

The regions of a tangle are shaded in a checkerboard fashion. With
\opt{sign=+} (or \opt{positive}, the default) the region at the star of the
outer disc, the leftmost region, is unshaded; \opt{sign=-} (or
\opt{negative}) shades it. \opt{shading=false} turns the shading off; if a
disc has an odd number of strands, no checkerboard shading exists and the
package warns and draws the tangle without it. Every disc gets a star in its
first region, and \opt{stars=false} leaves the stars out. \opt{shade} sets
the colour of the shaded regions.

\begin{code}
\PA{2,2}{x}  \PA[negative]{2,2}{x}  \PA[shading=false]{2,2}{x}
\PA[stars=false]{2,2}{x}  \PA[shade=red!20]{2,2}{x}
\end{code}
\begin{result}
  \PA{2,2}{x}\quad \PA[negative]{2,2}{x}\quad \PA[shading=false]{2,2}{x}\quad
  \PA[stars=false]{2,2}{x}\quad \PA[shade=red!20]{2,2}{x}
\end{result}

\subsection{Numbers and labels}

\opt{numbers} writes the number of every input disc into it (instead of its
label), in the order in which the discs appear in the layers. These are the
numbers used by \opt{insert} below. \opt{labels} numbers the boundary points
of the outer disc, clockwise from the star.

\begin{code}
\PA[numbers]{4,4}{D(2,2), A(1,1), id}\PA{4,4}{id, id, B(2,2)}\PA{4,2}{id, id, ev}
\PA[labels]{4,4}{D(2,2), A(1,1), id}\PA{4,4}{id, id, B(2,2)}\PA{4,2}{id, id, ev}
\end{code}
\begin{result}
  \PA[numbers]{4,4}{D(2,2), A(1,1), id}\PA{4,4}{id, id, B(2,2)}\PA{4,2}{id, id, ev}
  \qquad
  \PA[labels]{4,4}{D(2,2), A(1,1), id}\PA{4,4}{id, id, B(2,2)}\PA{4,2}{id, id, ev}
\end{result}

\subsection{Reduced form}

Two tangles describe the same element of a planar algebra when they are
isotopic, and a closed loop can be replaced by the scalar $\delta$, the loop
value. With \opt{reduce} the package computes a normal form: it removes
closed loops (they become a coefficient $\delta^k$ in front of the tangle),
pulls the strings tight so that zig-zags disappear and strands that only
turn around next to a disc are retracted, and draws the result in a standard
shape. In the normal form the input discs are stacked one above the other,
and between two consecutive discs (and above the first and below the last)
the strings form a Temperley--Lieb diagram written as cups, then through
strands, then caps. Discs on neighbouring levels can often be interchanged,
one passing the other on the left or on the right; the package goes through
all the arrangements of the discs that can be reached in this way and takes
the one whose description comes first in a fixed (lexicographic) order. So
isotopic tangles give the same normal form, and reducing a normal form gives
it back. Tangles with many discs that can move independently of each other
(for instance several discs that are not connected to the boundary) have
many arrangements; the search stops after \opt{max states} of them (default
\opt{100}) with a warning, and then the normal form may depend on how the
tangle was entered.

\begin{code}
\[ \PA{0,2}{coev}\PA{2,4}{id, coev, id}\PA{4,2}{id, ev, id}\PA{2,0}{ev}
   = \PA[reduce]{0,2}{coev}\PA{2,4}{id, coev, id}\PA{4,2}{id, ev, id}\PA{2,0}{ev} \]
\[ \PA{1,1}{id}\PA{1,3}{id, coev}\PA{3,1}{ev, id}
   = \PA[reduce]{1,1}{id}\PA{1,3}{id, coev}\PA{3,1}{ev, id}
   \qquad
   \PA{0,2}{coev}\PA{2,2}{T(2,2)}\PA{2,0}{ev}
   = \PA[reduce]{0,2}{coev}\PA{2,2}{T(2,2)}\PA{2,0}{ev} \]
\end{code}
\begin{result}
  \( \PA{0,2}{coev}\PA{2,4}{id, coev, id}\PA{4,2}{id, ev, id}\PA{2,0}{ev}
     = \PA[reduce]{0,2}{coev}\PA{2,4}{id, coev, id}\PA{4,2}{id, ev, id}\PA{2,0}{ev} \)
\end{result}
\begin{result}
  \( \PA{1,1}{id}\PA{1,3}{id, coev}\PA{3,1}{ev, id}
     = \PA[reduce]{1,1}{id}\PA{1,3}{id, coev}\PA{3,1}{ev, id}
     \qquad
     \PA{0,2}{coev}\PA{2,2}{T(2,2)}\PA{2,0}{ev}
     = \PA[reduce]{0,2}{coev}\PA{2,2}{T(2,2)}\PA{2,0}{ev} \)
\end{result}

\begin{code}
\[ \PA{2,4}{A(1,1), id, coev}\PA{4,2}{id, ev, B(1,1)}
   = \PA[reduce]{2,4}{A(1,1), id, coev}\PA{4,2}{id, ev, B(1,1)} \]
\end{code}
\begin{result}
  \( \PA{2,4}{A(1,1), id, coev}\PA{4,2}{id, ev, B(1,1)}
     = \PA[reduce]{2,4}{A(1,1), id, coev}\PA{4,2}{id, ev, B(1,1)} \)
\end{result}

The key \opt{loop value} sets the symbol of a loop (default \verb|\delta|),
and \opt{coefficient=false} draws the reduced tangle without the
coefficient.

\begin{code}
\PA[horizontal, reduce, loop value=n]{0,2}{coev}\PA{2,0}{ev}\PA{0,2}{coev}\PA{2,2}{x}\PA{2,0}{ev}
\end{code}
\begin{result}
  \PA[horizontal, reduce, loop value=n]{0,2}{coev}\PA{2,0}{ev}\PA{0,2}{coev}\PA{2,2}{x}\PA{2,0}{ev}
\end{result}

After every tangle (reduced or not) the following macros expand to data of
the tangle that was just drawn: \cs{PAlastloops} (the number of closed
loops), \cs{PAlastboxes} (the number of input discs) and, after
\opt{reduce}, \cs{PAlastnormalform} (the normal form as a string of
\cs{PA} commands, which can be fed back to the package). For the last
example, \cs{PAlastloops} is \PAlastloops, \cs{PAlastboxes} is
\PAlastboxes, and the normal form is
\begin{center}\ttfamily\small\PAlastnormalform\end{center}

\subsection{Defining and composing tangles}

\cs{PAdefine}\oarg{options}\marg{name}\marg{layers} stores a tangle under a
name instead of drawing it (with \opt{reduce}, its normal form is stored).
The key \opt{insert}\verb|={|$i$\verb|=|\meta{name}\verb|, ...}| glues the
named tangle into the input disc number~$i$ (as shown by \opt{numbers}); its
numbers of strands must agree with those of the disc. With \opt{expand},
every input disc whose label is the name of a tangle is replaced by it.

\begin{code}
\PAdefine{T}{\PA{2,4}{coev, id, id}\PA{4,2}{id, ev, S(1,1)}}
\[ \PA[numbers]{4,4}{D(2,2), A(1,1), id}\PA{4,4}{id, id, B(2,2)}\PA{4,2}{id, id, ev}
   \circ_3 \PA{2,4}{coev, id, id}\PA{4,2}{id, ev, S(1,1)}
   = \PA[insert={3=T}]{4,4}{D(2,2), A(1,1), id}\PA{4,4}{id, id, B(2,2)}\PA{4,2}{id, id, ev} \]
\[ \PA[insert={3=T}, reduce]{4,4}{D(2,2), A(1,1), id}\PA{4,4}{id, id, B(2,2)}\PA{4,2}{id, id, ev}
   \qquad
   \PA{2,2}{T}\PA{2,2}{T} = \PA[expand]{2,2}{T}\PA{2,2}{T} \]
\end{code}
\begin{result}
  \( \PA[numbers]{4,4}{D(2,2), A(1,1), id}\PA{4,4}{id, id, B(2,2)}\PA{4,2}{id, id, ev}
     \circ_3 \PA{2,4}{coev, id, id}\PA{4,2}{id, ev, S(1,1)}
     = \PA[insert={3=T}]{4,4}{D(2,2), A(1,1), id}\PA{4,4}{id, id, B(2,2)}\PA{4,2}{id, id, ev} \)
\end{result}
\begin{result}
  \( \PA[insert={3=T}, reduce]{4,4}{D(2,2), A(1,1), id}\PA{4,4}{id, id, B(2,2)}\PA{4,2}{id, id, ev}
     \qquad
     \PA{2,2}{T}\PA{2,2}{T} = \PA[expand]{2,2}{T}\PA{2,2}{T} \)
\end{result}

\section{Appearance}

\begin{center}
\begin{tabular}{@{}lll@{}}
  \toprule
  key & default & meaning \\
  \midrule
  \opt{strand sep}  & \opt{0.5}      & distance between neighbouring strands (cm) \\
  \opt{length}      & \opt{0.8}      & length of a layer (cm) \\
  \opt{box size}    & \opt{0.22}     & minimal radius of an input disc (cm) \\
  \opt{margin}      & \opt{0.3}      & space between the strings and the outer disc (cm) \\
  \opt{label font}  & \verb|\small|  & font of the disc labels \\
  \opt{scale}       & \opt{1}        & scales the picture (not the text) \\
  \opt{baseline}    & \opt{-0.5ex}   & vertical alignment, as in TikZ \\
  \opt{tikz}        &                & further options for the \texttt{tikzpicture} \\
  \opt{shade}       &                & colour of the shaded regions \\
  \opt{separator}   &                & symbol between the pictures of \cs{PAgenerators} \\
  \opt{gap}         & \opt{0.3}      & space between the pictures of \cs{PAgenerators} (cm) \\
  \bottomrule
\end{tabular}
\end{center}

The drawing uses the TikZ styles \opt{pa string}, \opt{pa shaded},
\opt{pa box} (input discs), \opt{pa outer}, \opt{pa star}, \opt{pa label}
(disc labels), \opt{pa number} (disc numbers), \opt{pa boundary label},
\opt{pa name}, \opt{pa separator}, \opt{pa coefficient} and \opt{pa error}.
The keys \opt{string style}, \opt{shade style}, \opt{box style},
\opt{outer style} and \opt{star style} append to the corresponding styles.

\begin{code}
\PA[string style={thick, blue!60!black}, outer style={line width=2pt}]{2,2}{x}
\PA[box style={fill=yellow}, star style={text=red}]{2,2}{x}
\PA[length=1.2, strand sep=0.8, box size=0.3]{3,3}{id, x(2,2)}\PA{3,1}{ev, id}
\end{code}
\begin{result}
  \PA[string style={thick, blue!60!black}, outer style={line width=2pt}]{2,2}{x}\qquad
  \PA[box style={fill=yellow}, star style={text=red}]{2,2}{x}\qquad
  \PA[length=1.2, strand sep=0.8, box size=0.3]{3,3}{id, x(2,2)}\PA{3,1}{ev, id}
\end{result}

\section{Layout}

The numbers of strands change from layer to layer, so their positions are
computed. Every item says where it would like the strands next to it to be:
a strand keeps its position, a cup or cap puts its two ends half a strand
distance on either side of its middle, and an input disc spreads its strands
evenly around its centre. On every seam between two layers the positions are
the best least-squares fit to these wishes under the condition that
neighbouring strands are at least \opt{strand sep} apart (an isotonic
regression), and a few sweeps back and forth through the tangle let both
neighbours of a seam have their say. The outer disc is then fitted around the
picture, and the strands that reach it are extended to its boundary.

\section{The standalone class}

With \verb|\documentclass[tikz]{standalone}| every TikZ picture becomes a
page of its own. Tangles inside \verb|\[...\]| and \verb|\(...\)| are handled
like in the \textsf{temperleylieb} and \textsf{plumber} packages: the whole
formula becomes one cropped page. Tangles outside formulas are pages of
their own. Other math environments (\verb|$...$|, \verb|equation|) cannot be
turned into pages; the package warns and draws the tangle inline, so use
\verb|\[...\]| there, or the class option \opt{varwidth}.

\section{Summary of commands}

\begin{description}
  \item[\cs{PA}\oarg{options}\marg{m,n}\marg{items}] One layer; a chain of
    layers is one tangle. \cs{PAlayer} is a synonym.
  \item[\cs{PAgenerators}\texttt{*}\oarg{options}] The items in \opt{list}
    side by side.
  \item[\cs{PAdefine}\oarg{options}\marg{name}\marg{layers}] Stores a
    tangle for \opt{insert} and \opt{expand}.
  \item[\cs{paset}\marg{options}] Sets options for the current group.
  \item[\cs{PAlastloops}, \cs{PAlastboxes}] Number of closed loops and of
    input discs of the last tangle.
  \item[\cs{PAlastnormalform}] Normal form of the last reduced tangle.
\end{description}

The package is written in \textsf{expl3} and needs a \LaTeX{} kernel from
2020-10 or newer and TikZ. It is distributed under the \LaTeX{} Project
Public License~1.3c.

\end{document}
